\section{Introduction}\label{sec:introduction}
Rewriting a binary changes the evidence available to an analyst without
necessarily changing the computation that the binary performs. This
distinction motivates the study of metamorphic transformations, both as
a software-protection mechanism and as a challenge for malware
analysis~\cite{PEsurveybible,metamorphic1,metamorphic2}. A useful assessment
must establish what was transformed, whether the intended behavior was
retained, and how the resulting representation was measured. These
questions become particularly difficult when executable instructions and
embedded data share a raw binary address space.

Zone provides a payload-level setting in which to investigate these
questions. Its binary-to-binary interface accepts compact Windows x86
and x64 inputs without source code or debugging information. The original
engine implements instruction substitution, insertion, and control-flow
transformations, including the Position-Independent Instruction (PII)
representation introduced in our SECRYPT 2026
paper~\cite{zone_secrypt_2026}. The present work retains that scope while
expanding the analysis and transformation paths.

\paragraph{Motivation.}
A sequence of bytes can be both a valid instruction encoding and part of
a string or another data object. Successful decoding therefore provides
insufficient evidence that the bytes should be rewritten as instructions.
The conference study identified this ambiguity as a source of functional
failures. Addressing it requires more than a different transformation:
the engine must identify candidate data, preserve its contents, and
maintain the references through subsequent layout changes. Failures at
any of these stages can invalidate an otherwise well-formed output.

LLMs introduce two further questions. They can be used to classify
regions of a binary, and they can be asked to transform an assembly
representation. The first task produces candidate labels or intervals;
the second produces a new instruction sequence. They have different
outputs, validation requirements, and failure modes. We consequently
study them as separate components, even though both interact with the
same data-handling infrastructure.

\paragraph{Relationship to the Conference Paper and Contributions.}
The original engine, PII design, and initial experiments appeared
in~\cite{zone_secrypt_2026}. The present article synthesizes that material
as a historical baseline and adds the following work:
\begin{enumerate}
\item A code--data identification layer combining bounded symbolic
exploration, static control-flow recovery, and complementary heuristics,
with an alternative LLM-based identification path and shared deterministic
data handling.
\item An LLM-assisted assembly-transformation pipeline for four
transformation families, with explicit zero-shot and two-shot conditions,
structural validation, and feedback after rejected attempts.
\item Integration of assembly wrapping, an existing technique, into Zone,
with the x64 implementation distinguished from the x86 port currently
under development.
\item An expanded assessment framework separating identification quality,
construction success, functional outcomes, static-scan observations,
similarity, diversity, and computational cost.
\end{enumerate}
The results of the additional campaigns are kept separate from the
published measurements. Reused visual material is attributed at its
point of use. The statement accompanying the submission details these
differences for the editors.

\paragraph{Scope.}
The subject is raw payload rewriting rather than complete Portable
Executable (PE) transformation. PE sections, imports, relocations, and
loader behavior introduce dependencies outside the experiments considered
here. Payload-level functional checks also capture only the behavior
observed by the chosen harness. Neither a valid assembly listing nor an
accepted raw-file scan establishes preservation of all program behaviors
or the response of a deployed endpoint product.

The article first introduces the relevant work
(Section~\ref{sec:background}) and Zone's architecture
(Section~\ref{sec:architecture}). Sections~\ref{sec:code_data},
\ref{sec:llm}, and~\ref{sec:wrapping} describe the extensions.
Section~\ref{sec:methodology} defines the assessment, and
Section~\ref{sec:results} separates historical and extended results.
Section~\ref{sec:discussion} discusses validity and remaining limitations;
Section~\ref{sec:conclusions} concludes.
